\chapter{2012年北京卷（理）}

\FigureLayoutDeclare{img/2012/beijing_science/q5_fig1.png}{5.6cm}{28bp 18bp 37bp 47bp}

\section{单选题}

\begin{problem}
已知集合 \(A = \{x \in \R \mid 3x + 2 > 0\}\)，\(B = \{x \in \R \mid (x + 1)(x - 3) > 0\}\)，则 \(A \cap B =\)（\quad）

\choices
    {\((-\infty, -1)\)}
    {\(\left(-1, -\frac{2}{3}\right)\)}
    {\(\left(-\frac{2}{3}, 3\right)\)}
    {\((3, +\infty)\)}
\end{problem}

\begin{answer}
D
\end{answer}

\begin{solution}
由 \(3x+2>0\) 得 \(x>-\frac{2}{3}\)，由 \((x+1)(x-3)>0\) 得 \(x<-1\) 或 \(x>3\). 因此

\[
A\cap B=\left(-\frac{2}{3},+\infty\right)\cap\left(\left(-\infty,-1\right)\cup(3,+\infty)\right)=(3,+\infty)
\]

故选 D.
\end{solution}

\begin{problem}
设不等式组 \(\begin{cases}
0 \leqslant x \leqslant 2,\\
0 \leqslant y \leqslant 2
\end{cases}\) 表示的平面区域为 \(D\). 在区域 \(D\) 内随机取一个点，则此点到坐标原点的距离大于 2 的概率是（\quad）

\choices
    {\(\frac{\pi}{4}\)}
    {\(\frac{\pi - 2}{2}\)}
    {\(\frac{\pi}{6}\)}
    {\(\frac{4 - \pi}{4}\)}
\end{problem}

\begin{answer}
D
\end{answer}

\begin{solution}
区域 \(D\) 是边长为 \(2\) 的正方形，面积为 \(4\). 在该正方形内，到原点距离不大于 \(2\) 的部分是半径为 \(2\) 的四分之一圆，面积为 \(\frac14\pi\cdot2^2=\pi\). 因而到原点距离大于 \(2\) 的部分面积为 \(4-\pi\)，所求概率为

\[
\frac{4-\pi}{4}
\]

故选 D.
\end{solution}

\begin{problem}
设 \(a\)，\(b \in \R\). “\(a=0\)”是“复数 \(a+b\i\) 是纯虚数”的（\quad）

\choices
    {充分而不必要条件}
    {必要而不充分条件}
    {充分必要条件}
    {既不充分也不必要条件}
\end{problem}

\begin{answer}
B
\end{answer}

\begin{solution}
复数 \(a+b\i\) 为纯虚数时，其实部必须为 \(0\)，且虚部不为 \(0\)，所以必有 \(a=0\). 反之，若 \(a=0\) 而 \(b=0\)，则 \(a+b\i=0\)，不满足纯虚数的非零虚部条件，因此 \(a=0\) 不是充分条件. 故“\(a=0\)”是“复数 \(a+b\i\) 是纯虚数”的必要而不充分条件，选 B.
\end{solution}

\begin{problem}
执行如图所示的程序框图，输出的 \(S\) 值为（\quad）

\texfigure[max width=0.70\linewidth]{img_repaint/2012/beijing_science/q4_fig1.tex}

\choices
    {\(2\)}
    {\(4\)}
    {\(8\)}
    {\(16\)}
\end{problem}

\begin{answer}
C
\end{answer}

\begin{solution}
初始时 \(k=0\)，\(S=1\). 每轮先判断 \(k<3\)，若成立则执行 \(S\leftarrow S\cdot 2^k\)，再执行 \(k\leftarrow k+1\)，然后返回判断框.

三次循环后的数值依次为 \((S,k)=(1,1)\)，\((S,k)=(2,2)\)，\((S,k)=(8,3)\).

此时 \(k<3\) 不成立，输出 \(S=8\). 故选 C.
\end{solution}

\begin{problem}
如图，\(\angle ACB = 90^{\circ}\)，\(CD \perp AB\) 于点 \(D\)，以 \(BD\) 为直径的圆与 \(BC\) 交于点 \(E\)，则（\quad）

\bitmapfigure[max width=0.42\linewidth]{img/2012/beijing_science/q5_fig1.png}

\choices
    {\(CE \cdot CB = AD \cdot DB\)}
    {\(CE\cdot CB = AD\cdot AB\)}
    {\(AD\cdot AB = CD^{2}\)}
    {\(CE\cdot EB = CD^{2}\)}
\end{problem}

\begin{answer}
A
\end{answer}

\begin{solution}
连接 \(DE\). 因为 \(BD\) 是圆的直径，所以 \(\angle BED=90^\circ\)，又 \(B\)，\(E\)，\(C\) 三点共线，故 \(DE\perp BC\). 在直角三角形 \(BCD\) 中，\(DE\) 是斜边 \(BC\) 上的高，所以

\[
CD^2=CE\cdot CB
\]

又因为 \(\triangle ABC\) 在 \(C\) 处直角，且 \(CD\perp AB\)，所以 \(\triangle ACD\sim\triangle CBD\)，从而

\[
CD^2=AD\cdot DB
\]

因此 \(CE\cdot CB=AD\cdot DB\)，故选 A.
\end{solution}

\begin{problem}
从 \(0\)，\(2\) 中选一个数字，从 \(1\)，\(3\)，\(5\) 中选两个数字，组成无重复数字的三位数，其中奇数的个数为（\quad）

\choices
    {\(24\)}
    {\(18\)}
    {\(12\)}
    {\(6\)}
\end{problem}

\begin{answer}
B
\end{answer}

\begin{solution}
若选出的偶数是 \(0\)，则 \(0\) 只能放在十位. 先从 \(1\)，\(3\)，\(5\) 中选两个数字，有 \(\mathrm C_3^2\) 种；个位可放其中任一个奇数，百位随之确定，所以有 \(\mathrm C_3^2\cdot2=6\) 个奇数.

若选出的偶数是 \(2\)，则从 \(1\)，\(3\)，\(5\) 中选两个数字后，个位有 \(2\) 种选法；剩下的 \(2\) 和另一个奇数可在百位、十位交换，有 \(2\) 种排法，所以有 \(\mathrm C_3^2\cdot2\cdot2=12\) 个奇数.

总数为 \(6+12=18\)，故选 B.
\end{solution}

\begin{problem}
某三棱锥的三视图如图所示，该三棱锥的表面积是（\quad）

\bitmapfigure[max width=0.76\linewidth]{img/2012/beijing_science/q7_fig1.png}

\choices
    {\(28 + 6\sqrt{5}\)}
    {\(30 + 6\sqrt{5}\)}
    {\(56 + 12\sqrt{5}\)}
    {\(60 + 12\sqrt{5}\)}
\end{problem}

\begin{answer}
B
\end{answer}

\begin{solution}
由三视图可还原为底面为直角三角形的三棱锥. 设底面顶点为 \(A=(0,0,0)\)，\(B=(5,0,0)\)，\(C=(5,4,0)\)，顶点为 \(P=(2,0,4)\). 这些坐标给出底面直角边长 \(5,4\)、顶点高 \(4\)，且顶点在底面边 \(AB\) 上距 \(A\) 为 \(2\) 的点的正投影.

底面面积为 \(S_{ABC}=\frac12\cdot5\cdot4=10\). 对于面 \(PBC\)，有 \(\overrightarrow{BP}=(-3,0,4)\)，\(\overrightarrow{BC}=(0,4,0)\)，故 \(\left\lvert\overrightarrow{BP}\times\overrightarrow{BC}\right\rvert=20\)，从而 \(S_{PBC}=10\). 对于面 \(PAB\)，有 \(\overrightarrow{AP}=(2,0,4)\)，\(\overrightarrow{AB}=(5,0,0)\)，故 \(\left\lvert\overrightarrow{AP}\times\overrightarrow{AB}\right\rvert=20\)，从而 \(S_{PAB}=10\).

对于面 \(PAC\)，有 \(\overrightarrow{AP}=(2,0,4)\)，\(\overrightarrow{AC}=(5,4,0)\).

故

\[
S_{PAC}=\frac12\left\lvert\overrightarrow{AP}\times\overrightarrow{AC}\right\rvert=\frac12\sqrt{(-16)^2+20^2+8^2}=6\sqrt5
\]

所以三棱锥的表面积为 \(10+10+10+6\sqrt5=30+6\sqrt5\)，故选 B.
\end{solution}

\begin{problem}
某棵果树前 \(n\) 年的总产量 \(S_n\) 与 \(n\) 之间的关系如图所示. 从目前记录的结果看，前 \(m\) 年的年平均产量最高，\(m\) 的值为（\quad）

\bitmapfigure[max width=0.5\linewidth]{img/2012/beijing_science/q8_fig1.png}

\choices
    {\(5\)}
    {\(7\)}
    {\(9\)}
    {\(11\)}
\end{problem}

\begin{answer}
C
\end{answer}

\begin{solution}
图中点 \((n,S_n)\) 表示前 \(n\) 年的总产量，因此前 \(n\) 年的年平均产量为 \(\frac{S_n}{n}\)，它等于原点与点 \((n,S_n)\) 的连线斜率. 比较图中各点与原点连线的斜率可知，斜率在 \(n=9\) 时最大，故前 \(9\) 年的年平均产量最高，选 C.
\end{solution}

\section{填空题}

\begin{problem}
直线 \(\begin{cases}
x=2+t,\\
y=-1-t,
\end{cases}\)（\(t\) 为参数）与曲线 \(\begin{cases}
x=3\cos\alpha,\\
y=3\sin\alpha,
\end{cases}\)（\(\alpha\) 为参数）的交点个数为\fillinblank{}.
\end{problem}

\begin{answer}
\(2\)
\end{answer}

\begin{solution}
直线参数方程给出 \(x+y=(2+t)+(-1-t)=1\)，即直线方程为 \(x+y=1\). 曲线参数方程给出 \(x^2+y^2=9\)，是以原点为圆心、半径为 \(3\) 的圆.

圆心到直线的距离为 \(\frac{1}{\sqrt2}<3\)，所以直线与圆相交于两个不同的点，交点个数为 \(2\).
\end{solution}

\begin{problem}
已知 \(\{a_n\}\) 为等差数列，\(S_n\) 为其前 \(n\) 项和，若 \(a_1=\frac{1}{2}\)，\(S_2=a_3\)，则 \(a_2=\)\fillinblank{}；\(S_n=\)\fillinblank{}.
\end{problem}

\begin{answer}
\(1\)；\(\frac{n(n+1)}{4}\)
\end{answer}

\begin{solution}
设等差数列的公差为 \(d\). 由 \(a_1=\frac12\) 得 \(a_2=\frac12+d\)，\(a_3=\frac12+2d\)，而 \(S_2=1+d\). 由 \(S_2=a_3\) 得

\[
1+d=\frac12+2d
\]

所以 \(d=\frac12\)，从而 \(a_2=1\). 又

\[
S_n=\frac n2\left(2a_1+(n-1)d\right)=\frac n2\left(1+\frac{n-1}{2}\right)=\frac{n(n+1)}4
\]
\end{solution}

\begin{problem}
在 \(\triangle ABC\) 中，若 \(a=2\)，\(b+c=7\)，\(\cos B=-\frac{1}{4}\)，则 \(b=\)\fillinblank{}.
\end{problem}

\begin{answer}
\(4\)
\end{answer}

\begin{solution}
由余弦定理，在角 \(B\) 对边 \(b\) 的关系中

\[
b^2=a^2+c^2-2ac\cos B=4+c^2+c=c^2+c+4
\]

又 \(b+c=7\)，即 \(b=7-c\). 代入得

\[
(7-c)^2=c^2+c+4
\]

即 \(49-14c+c^2=c^2+c+4\)，所以 \(15c=45\)，解得 \(c=3\)，从而 \(b=7-c=4\).
\end{solution}

\begin{problem}
在直角坐标系 \(xOy\) 中，直线 \(l\) 过抛物线 \(y^2 = 4x\) 的焦点 \(F\)，且与该抛物线相交于 \(A\)，\(B\) 两点，其中点 \(A\) 在 \(x\) 轴上方. 若直线 \(l\) 的倾斜角为 \(60^\circ\)，则 \(\triangle OAF\) 的面积为\fillinblank{}.
\end{problem}

\begin{answer}
\(\sqrt{3}\)
\end{answer}

\begin{solution}
抛物线 \(y^2=4x\) 的焦点为 \(F=(1,0)\). 直线 \(l\) 的倾斜角为 \(60^\circ\)，故斜率为 \(\sqrt3\)，其方程为 \(y=\sqrt3(x-1)\). 代入抛物线方程得

\[
3(x-1)^2=4x
\]

即 \(3x^2-10x+3=0\)，所以 \(x=3\) 或 \(x=\frac13\). 由于 \(A\) 在 \(x\) 轴上方，取 \(A=(3,2\sqrt3)\). 以 \(OF=1\) 为底，点 \(A\) 到 \(x\) 轴的距离为 \(2\sqrt3\)，于是

\[
S_{\triangle OAF}=\frac12\cdot1\cdot2\sqrt3=\sqrt3
\]
\end{solution}

\begin{problem}
已知正方形 \(ABCD\) 的边长为 1，点 \(E\) 是 \(AB\) 边上的动点，则 \(\overrightarrow{DE} \cdot \overrightarrow{CB}\) 的值为\fillinblank{}；\(\overrightarrow{DE} \cdot \overrightarrow{DC}\) 的最大值为\fillinblank{}.
\end{problem}

\begin{answer}
\(1\)；\(1\)
\end{answer}

\begin{solution}
建立直角坐标系，令 \(A=(0,0)\)，\(B=(1,0)\)，\(C=(1,1)\)，\(D=(0,1)\)，设 \(E=(t,0)\)，其中 \(0\leqslant t\leqslant1\).

有 \(\overrightarrow{DE}=(t,-1)\)，\(\overrightarrow{CB}=(0,-1)\)，\(\overrightarrow{DC}=(1,0)\).

因此 \(\overrightarrow{DE}\cdot\overrightarrow{CB}=1\)，\(\overrightarrow{DE}\cdot\overrightarrow{DC}=t\).

当 \(0\leqslant t\leqslant1\) 时，\(t\) 的最大值为 \(1\)，所以两空分别为 \(1,1\).
\end{solution}

\begin{problem}
已知 \(f(x)=m(x-2m)(x+m+3)\)，\(g(x)=2^x-2\). 若同时满足条件：

\begin{circlelist}
\item \(\forall x\in\R\)，\(f(x)<0\) 或 \(g(x)<0\)；
\item \(\exists x\in(-\infty,-4)\)，\(f(x)g(x)<0\)，
\end{circlelist}

则 \(m\) 的取值范围是\fillinblank{}.
\end{problem}

\begin{answer}
\((-4,-2)\)
\end{answer}

\begin{solution}
因为 \(g(x)=2^x-2\) 在 \(x<1\) 时为负，在 \(x=1\) 时为零，在 \(x>1\) 时为正. 条件（1）因此要求 \(f(x)<0\) 对一切 \(x\geqslant1\) 成立.

若 \(m\geqslant0\)，则 \(f\) 不可能在 \(x\geqslant1\) 上恒为负，所以 \(m<0\). 此时 \(f\) 的开口向下，两个零点为 \(2m\) 和 \(-m-3\). 要使 \(f(x)<0\) 对一切 \(x\geqslant1\) 成立，两个零点都必须小于 \(1\). 其中 \(2m<1\) 已由 \(m<0\) 保证，而 \(-m-3<1\) 给出 \(m>-4\). 故条件（1）等价于 \(-4<m<0\).

在 \(x<-4\) 时 \(g(x)<0\)，条件（2）要求存在 \(x<-4\) 使 \(f(x)>0\)，也就是 \(x\) 位于 \(f\) 的两个零点之间. 在 \(-4<m<0\) 中，\(-m-3>-3>-4\). 若 \(m\geqslant-2\)，则 \(2m\geqslant-4\)，两个零点都不小于 \(-4\)，不存在满足条件的 \(x\). 若 \(-4<m<-2\)，则 \(2m<-4<-m-3\)，区间 \((2m,-4)\) 中的任意 \(x\) 都满足 \(x<-4\) 且 \(f(x)>0\). 因而

\[
m\in(-4,-2)
\]
\end{solution}

\section{解答题}

\begin{problem}
已知函数 \(f(x)=\frac{(\sin x-\cos x)\sin 2x}{\sin x}\).

\begin{enumerate}
\item 求 \(f(x)\) 的定义域及最小正周期；
\item 求 \(f(x)\) 的单调递增区间.
\end{enumerate}
\end{problem}

\begin{solution}
\begin{enumerate}
\item 分母不能为零，所以定义域为
\[
D=\{x\in\R\mid \sin x\ne0\}=\{x\in\R\mid x\ne k\pi,\ k\in\Z\}
\]
在定义域上，
\[
f(x)=\frac{(\sin x-\cos x)\sin2x}{\sin x}=2\cos x(\sin x-\cos x)=\sin2x-\cos2x-1
\]
上式可写成 \(\sqrt2\sin\left(2x-\frac{\pi}{4}\right)-1\)，其所有正周期均为 \(k\pi\)（\(k\) 为正整数），故最小正周期候选为 \(\pi\). 另一方面，\(f\left(\frac{\pi}{4}\right)=0\)，\(f\left(\frac{3\pi}{4}\right)=-2\)，二者不相等，所以 \(\frac{\pi}{2}\) 不是周期. 由于定义域也以 \(\pi\) 为周期，故最小正周期为 \(\pi\).
\item 对化简后的函数求导，得
\[
f'(x)=2\cos2x+2\sin2x=2\sqrt2\sin\left(2x+\frac{\pi}{4}\right)
\]
当
\[
2k\pi<2x+\frac{\pi}{4}<(2k+1)\pi
\]
时，\(f'(x)>0\)，即
\[
k\pi-\frac{\pi}{8}<x<k\pi+\frac{3\pi}{8}
\]
但 \(x=k\pi\) 不在定义域内，因此
单调递增区间为 \(\left[k\pi-\frac{\pi}{8},k\pi\right)\) 和 \(\left(k\pi,k\pi+\frac{3\pi}{8}\right]\)，其中 \(k\in\Z\).
\end{enumerate}
\end{solution}

\begin{problem}
如图 1，在 \(\mathrm{Rt}\triangle ABC\) 中，\(\angle C=90^\circ\)，\(BC=3\)，\(AC=6\). \(D\)，\(E\) 分别是 \(AC\)，\(AB\) 上的点，且 \(DE\parallel BC\)，\(DE=2\)，将 \(\triangle ADE\) 沿 \(DE\) 折起到 \(\triangle A_1DE\) 的位置，使 \(A_1C \perp CD\)，如图 2.

\begin{enumerate}
\item 求证：\(A_1C \perp\) 平面 \(BCDE\)；
\item 若 \(M\) 是 \(A_1D\) 的中点，求 \(CM\) 与平面 \(A_1BE\) 所成角的大小；
\item 线段 \(BC\) 上是否存在点 \(P\)，使平面 \(A_1DP\) 与平面 \(A_1BE\) 垂直？说明理由.
\end{enumerate}

\bitmapfigure[max width=0.76\linewidth]{img/2012/beijing_science/q16_fig1.png}
\end{problem}

\begin{solution}
\begin{enumerate}
\item 因为 \(DE\parallel BC\)，所以 \(\triangle ADE\sim\triangle ABC\). 由
\[
\frac{DE}{BC}=\frac{2}{3}=\frac{AD}{AC}
\]
得 \(AD=4\)，\(CD=2\). 由 \(AC\perp BC\)、\(CD\parallel AC\) 和 \(DE\parallel BC\) 得 \(CD\perp DE\). 折叠保持 \(A_1D=AD\)，且原来 \(AD\perp DE\)，所以折叠后 \(A_1D\perp DE\). 故 \(\overrightarrow{CA_1}=\overrightarrow{CD}+\overrightarrow{DA_1}\) 也垂直于 \(DE\). 题设给出 \(A_1C\perp CD\)，所以 \(A_1C\) 垂直于平面 \(BCDE\) 内相交的两条直线 \(CD\)、\(DE\)，从而直线 \(A_1C\) 垂直于平面 \(BCDE\).
\item 建立空间直角坐标系，使 \(C=(0,0,0)\)，\(B=(3,0,0)\)，\(D=(0,2,0)\)，\(E=(2,2,0)\). 设 \(A_1=(u,v,w)\). 由 \(A_1C\perp CD\) 得 \(v=0\)，再由 \(A_1D=4\)、\(A_1E=AE=\sqrt{AD^2+DE^2}=2\sqrt5\)，得
\(u^2+w^2+4=16\)，\((u-2)^2+w^2+4=20\).
两式相减得 \(u=0\)，再得 \(w^2=12\). 取图示一侧 \(w=2\sqrt3\)，于是 \(A_1=(0,0,2\sqrt3)\)，\(M=(0,1,\sqrt3)\)，所以 \(\overrightarrow{CM}=(0,1,\sqrt3)\).

平面 \(A_1BE\) 的一个法向量为
\[
\bs{n}=(2\sqrt3,\sqrt3,3)
\]
因为 \(\bs{n}\cdot\overrightarrow{BA_1}=0\)，\(\bs{n}\cdot\overrightarrow{BE}=0\). 设 \(CM\) 与平面 \(A_1BE\) 所成的角为 \(\theta\)，则
\[
\sin\theta=\frac{\left|\overrightarrow{CM}\cdot\bs{n}\right|}{|\overrightarrow{CM}|\,|\bs{n}|}
\]
其中 \(\left|\overrightarrow{CM}\right|=2\)，\(\left|\bs{n}\right|=2\sqrt6\)，\(\left|\overrightarrow{CM}\cdot\bs{n}\right|=4\sqrt3\). 因而 \(\sin\theta=\frac{\sqrt2}{2}\)，且 \(0\leqslant\theta\leqslant\frac{\pi}{2}\)，所以 \(\theta=45^\circ\).
\item 设 \(P=(p,0,0)\)，其中 \(0\leqslant p\leqslant3\). 平面 \(A_1DP\) 的一个法向量可取为
\[
\bs{n}_1=(2\sqrt3,\sqrt3p,p)
\]
若两平面垂直，则其法向量应垂直，即
\[
\bs{n}_1\cdot\bs{n}
=(2\sqrt3,\sqrt3p,p)\cdot(2\sqrt3,\sqrt3,3)
=12+6p=0
\]
这给出 \(p=-2\)，与 \(0\leqslant p\leqslant3\) 矛盾，因此线段 \(BC\) 上不存在满足条件的点 \(P\).
\end{enumerate}
\end{solution}

\begin{problem}
近年来，某市为促进生活垃圾的分类处理，将生活垃圾分为厨余垃圾，可回收物和其他垃圾三类，并分别设置了相应的垃圾箱. 为调查居民生活垃圾分类投放情况，现随机抽取了该市三类垃圾箱中总计 \(1000\text{ 吨}\) 生活垃圾，数据统计如下（单位：\(\text{吨}\)）：

\begin{center}
\renewcommand{\arraystretch}{1.2}
\begin{tabular}{|c|c|c|c|}
\hline
& “厨余垃圾”箱 & “可回收物”箱 & “其他垃圾”箱 \\
\hline
厨余垃圾 & 400 & 100 & 100 \\
\hline
可回收物 & 30 & 240 & 30 \\
\hline
其他垃圾 & 20 & 20 & 60 \\
\hline
\end{tabular}
\end{center}

\begin{enumerate}
\item 试估计厨余垃圾投放正确的概率；
\item 试估计生活垃圾投放错误的概率；
\item 假设厨余垃圾在“厨余垃圾”箱，“可回收物”箱，“其他垃圾”箱的投放量分别为 \(a\)，\(b\)，\(c\)，其中 \(a>0\)，\(a+b+c=600\). 当数据 \(a\)，\(b\)，\(c\) 的方差 \(s^2\) 最大时，写出 \(a\)，\(b\)，\(c\) 的值（结论不要求证明），并求此时 \(s^2\) 的值.
\end{enumerate}

（注：\(s^2=\frac{1}{n}\left[(x_1-\overline{x})^2+(x_2-\overline{x})^2+\cdots+(x_n-\overline{x})^2\right]\)，其中 \(\overline{x}\) 为数据 \(x_1\)，\(x_2\)，\(\cdots\)，\(x_n\) 的平均数）
\end{problem}

\begin{solution}
\begin{enumerate}
\item 厨余垃圾总量为 \(400+100+100=600\) 吨，其中投放正确的有 \(400\) 吨，所以厨余垃圾投放正确的概率估计为
\[
\frac{400}{600}=\frac{2}{3}
\]
\item 投放错误的垃圾总量为非主对角线数据之和
\[
100+100+30+30+20+20=300\text{ 吨}
\]
因此生活垃圾投放错误的概率估计为
\[
\frac{300}{1000}=\frac{3}{10}
\]
\item 因为 \(a+b+c=600\)，所以 \(a\)，\(b\)，\(c\) 的平均数为 \(200\)，从而
\[
s^2=\frac13\left[(a-200)^2+(b-200)^2+(c-200)^2\right]
\]
由于 \(b\)，\(c\geqslant0\)，有 \(b^2+c^2\leqslant(b+c)^2=(600-a)^2\)，等号当且仅当 \(b=0\) 或 \(c=0\). 又 \(0<a\leqslant600\)，所以
\[
a^2+b^2+c^2\leqslant a^2+(600-a)^2=2(a-300)^2+180000\leqslant360000
\]
最后一个等号在 \(a=600\) 时取得，此时 \(b=c=0\). 因此
\[
s^2=\frac13\left[(600-200)^2+(0-200)^2+(0-200)^2\right]=80000
\]
\end{enumerate}
\end{solution}

\begin{problem}
已知函数 \(f(x)=ax^2+1\)（\(a>0\)），\(g(x)=x^3+bx\).

\begin{enumerate}
\item 若曲线 \(y=f(x)\) 与曲线 \(y=g(x)\) 在它们的交点 \((1,c)\) 处具有公共切线，求 \(a\)，\(b\) 的值；
\item 当 \(a^2=4b\) 时，求函数 \(f(x)+g(x)\) 的单调区间，并求其在区间 \((-\infty,-1]\) 上的最大值.
\end{enumerate}
\end{problem}

\begin{solution}
\begin{enumerate}
\item 由交点 \((1,c)\) 同时在两条曲线上，得
\[
a+1=1+b
\]
所以 \(a=b\). 又因为两曲线在该点有公共切线，所以切线斜率相等，即
\[
2a=3+b
\]
联立 \(a=b\) 得 \(a=3\)，\(b=3\).
\item 令 \(h(x)=f(x)+g(x)=x^3+ax^2+\frac{a^2}{4}x+1\). 则
\[
h'(x)=3x^2+2ax+\frac{a^2}{4}=3\left(x+\frac a2\right)\left(x+\frac a6\right)
\]
由于 \(a>0\)，故 \(h'(x)>0\) 在 \(\left(-\infty,-\frac a2\right)\) 和 \(\left(-\frac a6,+\infty\right)\) 上成立，\(h'(x)<0\) 在 \(\left(-\frac a2,-\frac a6\right)\) 上成立. 因此 \(h\) 在 \(\left(-\infty,-\frac a2\right]\) 和 \(\left[-\frac a6,+\infty\right)\) 上单调递增，在 \(\left[-\frac a2,-\frac a6\right]\) 上单调递减.

又 \(h\left(-\frac a2\right)=1\)，且 \(h(-1)=a-\frac{a^2}{4}=1-\frac{(a-2)^2}{4}\).
当 \(0<a\leqslant2\) 时，\(-1\leqslant-\frac a2\)，函数在 \((-\infty,-1]\) 上递增，所以最大值为 \(h(-1)=a-\frac{a^2}{4}\). 当 \(a>2\) 时，点 \(-\frac a2\) 属于 \((-\infty,-1]\)，且 \(h(-1)\leqslant1\)，故该区间上的最大值为 \(h\left(-\frac a2\right)=1\).
\end{enumerate}
\end{solution}

\begin{problem}
已知曲线 \(C:(5-m)x^2+(m-2)y^2=8\)（\(m\in\R\)）.

\begin{enumerate}
\item 若曲线 \(C\) 是焦点在 \(x\) 轴上的椭圆，求 \(m\) 的取值范围；
\item 设 \(m=4\)，曲线 \(C\) 与 \(y\) 轴的交点为 \(A\)，\(B\)（点 \(A\) 位于点 \(B\) 的上方），直线 \(y=kx+4\) 与曲线 \(C\) 交于不同的两点 \(M\)，\(N\)，直线 \(y=1\) 与直线 \(BM\) 交于点 \(G\). 求证：\(A\)，\(G\)，\(N\) 三点共线.
\end{enumerate}
\end{problem}

\begin{solution}
\begin{enumerate}
\item 要使 \(C\) 为椭圆，必须有 \(5-m>0\) 且 \(m-2>0\)，即 \(2<m<5\). 标准化得
\[
\frac{x^2}{8/(5-m)}+\frac{y^2}{8/(m-2)}=1
\]
焦点在 \(x\) 轴上要求 \(x\) 方向的半长轴较长，即
\[
\frac{8}{5-m}>\frac{8}{m-2}
\]
从而 \(m>\frac72\). 所以 \(m\in\left(\frac72,5\right)\).
\item 当 \(m=4\) 时，曲线为 \(x^2+2y^2=8\)，故 \(A=(0,2)\)，\(B=(0,-2)\).
设 \(M=(u,ku+4)\)，\(N=(v,kv+4)\).
把直线 \(y=kx+4\) 代入曲线，得
\[
(1+2k^2)x^2+16kx+24=0
\]
由韦达定理，\(u+v=-\frac{16k}{1+2k^2}\)，\(uv=\frac{24}{1+2k^2}\)，所以
\[
2kuv+3(u+v)=\frac{48k}{1+2k^2}-\frac{48k}{1+2k^2}=0
\]

直线 \(BM\) 与 \(y=1\) 交于 \(G\). 由 \(B=(0,-2)\) 和 \(M=(u,ku+4)\)，可得
\[
x_G=\frac{3u}{ku+6}
\]
另一方面，直线 \(AN\) 与 \(y=1\) 的交点横坐标为
\[
-\frac{v}{kv+2}
\]
而
\[
3u(kv+2)+v(ku+6)=2\bigl(2kuv+3(u+v)\bigr)=0
\]
故
\[
\frac{3u}{ku+6}=-\frac{v}{kv+2}
\]
因此 \(G\) 同时在直线 \(AN\) 和直线 \(y=1\) 上，故 \(A\)，\(G\)，\(N\) 三点共线.
\end{enumerate}
\end{solution}

\begin{problem}
设 \(A\) 是由 \(m\times n\) 个实数组成的 \(m\) 行 \(n\) 列的数表，满足：每个数的绝对值不大于 1，且所有数的和为零. 记 \(S(m,n)\) 为所有这样的数表构成的集合. 对于 \(A\in S(m,n)\)，记 \(r_i(A)\) 为 \(A\) 的第 \(i\) 行各数之和（\(1\le i\le m\)），\(c_j(A)\) 为 \(A\) 的第 \(j\) 列各数之和（\(1\le j\le n\)）；记 \(k(A)\) 为 \(|r_1(A)|\)，\(|r_2(A)|\)，\(\cdots\)，\(|r_m(A)|\)，\(|c_1(A)|\)，\(|c_2(A)|\)，\(\cdots\)，\(|c_n(A)|\) 中的最小值.

\begin{enumerate}
\item 对如下数表 \(A\)，求 \(k(A)\) 的值：

\begin{center}
\renewcommand{\arraystretch}{1.2}
\begin{tabular}{|c|c|c|}
\hline
1 & 1 & \(-0.8\) \\
\hline
\(0.1\) & \(-0.3\) & \(-1\) \\
\hline
\end{tabular}
\end{center}

\item 设数表 \(A\in S(2,3)\) 形如

\begin{center}
\renewcommand{\arraystretch}{1.2}
\begin{tabular}{|c|c|c|}
\hline
1 & 1 & \(c\) \\
\hline
\(a\) & \(b\) & \(-1\) \\
\hline
\end{tabular}
\end{center}

求 \(k(A)\) 的最大值；
\item 给定正整数 \(t\)，对于所有的 \(A\in S(2,2t+1)\)，求 \(k(A)\) 的最大值.
\end{enumerate}
\end{problem}

\begin{solution}
\begin{enumerate}
\item 两行的和分别为 \(1+1-0.8=1.2\) 和 \(0.1-0.3-1=-1.2\)，三列的和分别为 \(1.1\)，\(0.7\)，\(-1.8\). 所有行和、列和的绝对值依次为 \(1.2\)，\(1.2\)，\(1.1\)，\(0.7\)，\(1.8\)，故
\[
k(A)=0.7
\]
\item 由所有元素的和为零，得 \(1+1+c+a+b-1=0\)，从而 \(a+b+c=-1\).
因此两行和分别为 \(2+c\) 和 \(-(2+c)\)，三列和分别为 \(1+a\)，\(1+b\)，\(c-1\). 由于 \(a\)，\(b\)，\(c\in[-1,1]\)，有
\[
k(A)=\min\{2+c,1+a,1+b,1-c\}
\]
若 \(k(A)>1\)，则 \(2+c>1\)，\(1+a>1\)，\(1+b>1\)，\(1-c>1\)，从而 \(c>-1\)，\(a>0\)，\(b>0\)，且 \(c<0\). 但 \(a+b=-1-c<0\)，矛盾，所以 \(k(A)\leqslant1\). 取 \(a=b=0\)，\(c=-1\) 时，矩阵满足条件且五个相关量为 \(1\)，\(1\)，\(1\)，\(1\)，\(2\)，故最大值为 \(1\).
\item 令 \(N=2t+1\)，并设 \(K=k(A)\). 由于所有元素的总和为零，设第一行和为 \(r\)，则第二行和为 \(-r\)，且 \(|r|\geqslant K\). 每列的列和记为 \(s_j\)，则 \(\sum_{j=1}^{N}s_j=0\)，且 \(|s_j|\geqslant K\).

若 \(K=0\)，结论立即成立. 以下设 \(K>0\). 由于 \(2t+1\) 为奇数，正列或负列至少有 \(t+1\) 个. 交换两行可使 \(r\geqslant K\)；若此时至少 \(t+1\) 个列和为负，则把所有元素同时取相反数并再次交换两行，仍有 \(r\geqslant K\)，且至少 \(t+1\) 个列和为正. 记这些正列的集合为 \(P\)，其余列为 \(Q\)，设 \(|P|=p\geqslant t+1\)，并令
\[
U=\sum_{j\in P}s_j=-\sum_{j\in Q}s_j
\]
对 \(j\in P\)，第一行元素 \(x_j\leqslant1\)；对 \(j\in Q\)，若第二行元素为 \(y_j\)，则 \(x_j=s_j-y_j\leqslant s_j+1\). 因此
\[
r=\sum_{j=1}^{N}x_j\leqslant p+\sum_{j\in Q}(s_j+1)=N-U
\]
又 \(U\geqslant pK\geqslant(t+1)K\)，所以
\[
K\leqslant r\leqslant N-(t+1)K
\]
从而
\[
K\leqslant\frac{N}{t+2}=\frac{2t+1}{t+2}
\]

下面构造达到该上界的数表. 令 \(K_0=\frac{2t+1}{t+2}\)，\(u=\frac{(t+1)K_0}{t}\). 当 \(t\geqslant1\) 时，\(1\leqslant K_0<2\)，且 \(K_0<u\leqslant2\).
对前 \(t+1\) 列取两行元素分别为 \(1\)、\(K_0-1\)，对后 \(t\) 列取两行元素分别为 \(1-u\)、\(-1\). 因为 \(K_0-1\in[0,1)\)、\(1-u\in[-1,0)\)，所以这些元素的绝对值均不大于 \(1\). 两行和分别为
\((t+1)+t(1-u)=K_0\) 和 \((t+1)(K_0-1)-t=-K_0\)，
前 \(t+1\) 列的列和为 \(K_0\)，后 \(t\) 列的列和为 \(-u\). 因而各行和的绝对值为 \(K_0\)，各列和的绝对值为 \(K_0\) 或 \(u\)，且 \(u>K_0\)，所以该数表的 \(k(A)\) 为 \(K_0\)，上界可以达到. 所以最大值为
\[
\frac{2t+1}{t+2}
\]
\end{enumerate}
\end{solution}
